\section{讨论与局限}
现有结果支持轨迹条件化的信息组织及共享训练下的本站部署，但证据来自不同协议：CBH-SL 支撑本站部署，条件与天气消融主要来自区域版本。本站 E2N/E2C 及其匹配消融尚未完成；当前归档也未恢复 SL 最终 checkpoint，完整逐样本复评仍受限。因此不能将这些结果拼接为已验证的本站增强模型。

后续需补齐本站增强结构与强基线的共同精度场景复评；迁移应计入 checkpoint 选择标签，或在开发站冻结训练步数。单一区域、缺测窗口、已查看年份与队列筛选限制外推；场站区间未覆盖全部空间依赖和模型选择，探索性分组也未统一校正。新确认应使用独立年份或站组。

\section{结论}
历史轨迹条件化天气修正在保留完整历史记忆的同时，显式提供待修正的未来功率判断。全 56 站结果支持共享训练、本站部署；匹配对照表明轨迹信息包的收益大于额外摘要及单纯输出相加的增量。区域非线性天气扩展进一步降低源站误差，但未见站适配未稳定超过强预训练 Transformer。区分训练来源、在线输入和标签预算，是解释这些收益及其适用范围的关键。
